\documentclass[10pt,a4paper]{amsart}
\usepackage[margin=19mm]{geometry}
\usepackage{amsmath,amssymb,mathtools}
\usepackage[scr=rsfs,cal=euler]{mathalfa}
\usepackage{xfrac}
\usepackage[unicode,hidelinks,bookmarks=false]{hyperref}
\newcommand{\ie}{i.e.\ }
\newcommand{\cf}{cf.\ }
\newcommand{\resp}{respectively\ }
\newcommand{\Subsection}{\subsection}
\providecommand{\nobreakdash}{\nobreak\hbox{-}}
\setlength{\emergencystretch}{2em}
\multlinegap0pt
\usepackage{enumitem}
\usepackage{breqn}
\usepackage{setspace}
\usepackage{bm}
\usepackage{mathtools}
\usepackage{cleveref}
\mathtoolsset{showonlyrefs}
\newcommand{\ba}{\bar{a}}
\newcommand{\bb}{\bar{b}}
\newcommand{\bi}{\bar{i}}
\newcommand{\bj}{\bar{j}}
\newcommand{\bk}{\bar{k}}
\newcommand{\bl}{\bar{l}}
\newcommand{\bn}{\bar{n}}
\newcommand{\bp}{\bar{p}}
\newcommand{\bq}{\bar{q}}
\newcommand{\br}{\bar{r}}
\newcommand{\bs}{\bar{s}}
\newcommand{\bt}{\bar{t}}
\newcommand{\bw}{\bar{w}}
\newcommand{\bz}{\bar{z}}
\newcommand{\bA}{\bar{A}}
\newcommand{\bB}{\bar{B}}
\newcommand{\bI}{\bar{I}}
\newcommand{\bJ}{\bar{J}}
\newcommand{\bK}{\bar{K}}
\newcommand{\bL}{\bar{L}}
\newcommand{\bM}{\bar{M}}
\newcommand{\bN}{\bar{N}}
\newcommand{\bP}{\bar{P}}
\newcommand{\bQ}{\bar{Q}}
\newcommand{\bR}{\bar{R}}
\newcommand{\bS}{\bar{S}}
\newcommand{\bT}{\bar{T}}
\newcommand{\bW}{\bar{W}}
\newcommand{\bZ}{\bar{Z}}
\newcommand{\p}{\phi}
\newcommand{\vn}{{\bf n}}
\newcommand{\balpha}{\bar{\alpha}}
\newcommand{\bbeta}{\bar{\beta}}
\newcommand{\eeta}{\bar{\eta}}
\newcommand{\btau}{\bar{\tau}}
\newcommand{\bxi}{\bar{\xi}}
\newcommand{\bzeta}{\bar{\zeta}}
\newcommand{\sbpartial}{\bar{\partial}^*}
\newcommand{\spartial}{\partial^*}
\newcommand{\bpartial}{\bar{\partial}}
\newcommand{\bnabla}{\overline{\nabla}}
\newcommand{\pbp}{\partial \bar{\partial}}
\newcommand{\bpp}{\bar{\partial}\partial}
\newcommand{\red}[1]{\textcolor{red}{#1}}
\newcommand{\tr}[2]{\mbox{tr}_{#1}{#2}}
\newcommand{\trw}[1]{\mbox{tr}_{\omega}{#1}}
\def \tg{\tilde{g}}
\def \bara{\bar{\alpha}}
\def \barb{\bar{\beta}}
\newcommand{\fg}{\mathfrak{g}}
\newcommand{\fI}{\mathfrak{I}}
\newcommand{\fIm}{\mathfrak{Im}}
\newcommand{\fRe}{\mathfrak{Re}}
\newcommand{\fRic}{\mathfrak{Ric}}
\newcommand{\bfC}{\hbox{\bbbld C}}
\newcommand{\bfH}{\hbox{\bbbld H}}
\newcommand{\bfN}{\hbox{\bbbld N}}
\newcommand{\bfP}{\hbox{\bbbld P}}
\newcommand{\bfQ}{\hbox{\bbbld Q}}
\newcommand{\bfR}{\hbox{\bbbld R}}
\newcommand{\bfS}{\hbox{\bbbld S}}
\newcommand{\bfT}{\hbox{\bbbld T}}
\newcommand{\bfZ}{\hbox{\bbbld Z}}
\newcommand{\cA}{\mathcal{A}}
\newcommand{\cB}{\mathcal{B}}
\newcommand{\cC}{\mathcal{C}}
\newcommand{\cF}{\mathcal{F}}
\newcommand{\cH}{\mathcal{H}}
\newcommand{\cK}{\mathcal{K}}
\newcommand{\cL}{\mathcal{L}}
\newcommand{\cM}{\mathcal{M}}
\newcommand{\cP}{\mathcal{P}}
\newcommand{\cT}{\mathcal{T}}
\newcommand{\Hess}{\mbox{Hess}}
\newcommand{\Rc}{\mbox{Rc}}
\newcommand{\Ric}{\mbox{Ric}}
\newcommand{\Ker}{\mbox{Ker}}
\newcommand{\tF}{\tilde{F}}
\newcommand{\tU}{\tilde{U}}
\newcommand{\va}{{\bf a}}
\newcommand{\ve}{{\bf e}}
\newcommand{\vv}{{\bf v}}
\newcommand{\ol}{\overline}
\newcommand{\ul}{\underline}
\newtheorem{theorem}{Theorem}[section]
\newtheorem{claim}{Claim}[section]
\newtheorem{lemma}[theorem]{Lemma}
\newtheorem{proposition}[theorem]{Proposition}
\newtheorem{corollary}[theorem]{Corollary}
\newtheorem{conjecture}[theorem]{Conjecture}
\newtheorem{definition}[theorem]{Definition}
\newtheorem{example}[theorem]{Example}
\newtheorem{problem}[theorem]{Problem}
\newtheorem{xca}[theorem]{Exercise}
\newtheorem{remark}[theorem]{Remark}
\numberwithin{equation}{section}
\newcommand{\abs}[1]{\lvert#1\rvert}
\begin{document}
\title[Monge-Amp`ere-type equation for forms of positive degree and]{Monge-Amp\`ere-type equation for forms of positive degree and Demailly's transcendental Morse inequality}
\author{Mathew George}
\date{}
\maketitle
\noindent\emph{Source and licence.} Monge-Amp\`ere-type equation for forms of positive degree and Demailly's transcendental Morse inequality. Version arXiv:2606.23981v2 (CC BY 4.0). This material is distributed under the Creative Commons Attribution 4.0 International licence, http://creativecommons.org/licenses/by/4.0/, as recorded in the arXiv OAI licence metadata for this exact version and shown on the abstract page https://arxiv.org/abs/2606.23981v2. Full rights notice: licenses/arxiv-2606.23981-CC-BY-4.0.txt.\par\smallskip
\noindent\textit{Editorial scope.} Counted unit: the original \texttt{theorem} environment \texttt{demailly-thm} at source lines 1018--1040 together with its complete proof at lines 1042--1132; the delivered document keeps source lines 258--1133 so that every referenced label and every citation resolves inside it. Native editable source is the paper's own concatenated TeX; the AMS wrapper is editorial formatting only. No publisher class, upstream script or unrelated artwork is redistributed.
\par\smallskip\noindent\textit{Reference note.} This article ships no compiled bibliography (biblatex); the entries delivered here are composed from the author's own BibTeX (.bib) fields, with no field invented and no entry dropped.
\begin{theorem}\label{exist-theorem}
  For any $(a,b)$ with $a+nb = m-n-2$ satisfying condition \ref{a} from Theorem \ref{ab-MAeqn-theorem}, there exists a form $u \in C^{\infty}_{m-1,m-1}(M, \mathbb R)$ and  a unique constant $c > 0$ satisfying

$$\left[\star_{\omega}\left((\alpha + i \pbp u)\wedge \omega_{n-m-1}\right)\right]^n = c \cdot dV,$$

\noindent with  

$$   \alpha + i \pbp u >0, \quad u \in \Ker (\spartial) \cap \Ker( \bar{\partial}^*),$$

\noindent and

$$ (m-1)\Lambda^{m-2} \Delta u = (1+a) i \pbp \Lambda^{m-1}  u +  (1+b) \Lambda^{m-1} \Delta u \;\omega.$$

\

If $(a,b)$ and $\alpha$ satisfies condition \ref{b} in Theorem \ref{ab-MAeqn-theorem}, with $X = \star_{\omega}(\alpha \wedge \omega_{n-m-1})$, then there exists a constant $\epsilon>0$ such that if $|dV - \log\det( \star_{\omega}(\alpha \wedge \omega_{n-m-1})) \omega^n|_{C^{0,\gamma}}< \epsilon$ for some $\gamma >0 $, then for some $c>0$ there exists a solution $u$ for this system.

\end{theorem}


\


The solution $u$ cannot be guaranteed to be unique, although its lowest order primitive coordinates $\psi_0$ and $\psi_1$ are unique up to harmonic forms. Further uniqueness of higher order primitives could only be established under additional assumptions as in Theorem \ref{unique-thm}. As a consequence of the existence of such a form $u \in C^{\infty}_{m-1,m-1}(M, \mathbb R)$ for $b=0$ and $a = m-n-2 $, the Demailly's transcendental Morse inequality for higher cohomologies (Theorem \ref{demailly-thm}) can now be proven.


Independent of the geometric application, we study this equation for other parametric values $(a,b)$. In Section \ref{section3}, this system is reduced to the {\em $(a,b)$ Monge-Amp\`ere-type equation}: 

$$    \left(\star_{\omega}(\alpha \wedge \omega_{n-m-1}) + \frac{n!}{(n-m+1)!} \left[ a i\pbp f + b \Delta f \omega\right] \right)^n = dV.$$

\

We prove the following.

\begin{theorem} \label{ab-MAeqn-theorem}
For any pair $(a,b) \neq (0,0)$, consider the equation

\begin{equation}\label{ab-MAeqn1}
    \log\det(\tilde g_f) = c.\psi(z), \text{ with }\;\;\; \tilde g_f >0.
\end{equation}

\noindent where

$$\tilde g_f = X + \frac{n!}{(n-m+1)!} \left[ a i\pbp f + b \Delta f \omega\right],$$

\

\noindent with $X$ a positive $(1,1)$ form and $\psi(z) \in C^{\infty}(M)$ a positive function. Then the following statements are true.

\begin{enumerate}[label=\Roman*.]
    \item \label{a} If $a+b\geq 0$ and $b \geq 0$, \textbf{or} if $a+b \leq 0$ and $b \leq 0$, then for some $c>0$, there exists a solution $f$ for \eqref{ab-MAeqn1} unique up to constants.

    
    \item \label{b} If $a+b < 0$, $b > 0$ \textbf{or} $a+b >0$, $b < 0$, and $\kappa(X) < \dfrac{|a|}{n|b|}$, then there exists a constant $\epsilon>0$ such that if $|\psi - \log\det(X)|_{C^{0, \gamma}}< \epsilon$ for some $\gamma >0 $, then for some $c>0$ there exists a solution $f$ for \eqref{ab-MAeqn1} unique up to constants. Here $\kappa(X)$ denotes the condition number of $X$.
\end{enumerate}





\end{theorem}



\

In Section $2$, we establish the strengthened uniqueness theorem under the Kernel condition and also prove that the original Laplace Trace condition is too rigid. In Section $3$, we introduce the $(a,b)$ Monge-Amp\`ere-type equation, and show the application to analytical Demailly's Morse inequality. Finally in Section $4$, \emph{a priori} estimates are shown which implies Theorem \ref{exist-theorem} and Theorem \ref{ab-MAeqn-theorem}.


\section*{Acknowledgement and AI tool use}

The author would like to thank Sławomir Dinew for helpful communications. 

Gemini was used for literature search and improved writing at a few places.



\section{The Kernel Condition and Uniqueness}

\bigskip


The following notations will be used in this paper. For any $(m-1,m-1)$ form $\alpha$ and any $(1,1)$ form $\beta$

\begin{itemize}

    \item $\alpha_{\omega} = \star_{\omega}(\alpha \wedge \omega_{n-m-1})$ defines a $(1,1)$ form.

    \

    \item $\beta_{k}=\dfrac{\beta^k}{k!}.$

    \

    \item The $\partial$-Laplacian of the K\"ahler metric $\omega$ will simply be written $\Delta$ which is defined by $\partial \spartial + \spartial \partial$.

    \

    \item $\Lambda$ and $L$ denote the Trace operator and the Lefschetz operator for the metric $\omega$ respectively.
\end{itemize}



We recall the definition of strong positivity.

\begin{definition}
    An $(m,m)$ form $\alpha \in C^{\infty}_{m,m}(M,\mathbb R)$ is called {\bf strongly positive} if it can be decomposed as 

     $$\alpha =  \sum_{s} c_s \, (i\alpha^s_1 \wedge \bar{\alpha^s}_1) \wedge \dots \wedge (i\alpha^s_m \wedge \bar{\alpha^s}_m)$$

     \noindent for $c_s >0$ and $\alpha^s_i$ a $(1,0)$ form. We denote this by $\alpha >0$. If $c_s \geq 0$, then $\alpha$ is strongly semi-positive, denoted by $\alpha \geq 0$. 
\end{definition}






Given an $\alpha \in C^{\infty}_{m,m}(M, \mathbb R)$ with $\alpha >0$, the Monge-Amp\`ere equation for forms of positive degree for $u \in C_{m-1,m-1}^{\infty}(M, \mathbb R)$ can be written as

\begin{equation}\label{MA-eqn}
    \left[\star_{\omega}\left((\alpha + i \pbp u)\wedge \omega_{n-m-1}\right)\right]^n = dV
\end{equation}

\ 

\noindent assuming $\alpha + i \pbp u >0$.


\bigskip


{\flushleft \bm{{\bf The Kernel Condition $u \in \text{Ker}(\partial^*) \cap \text{Ker}(\bar{\partial}^*)$:}} }
\
\bigskip 

 
Let $u \in C^{\infty}_{m-1,m-1}(M, \mathbb R)$ be a real form in $\mbox{Ker}(\partial^*)\cap \text{Ker}(\bar{\partial}^*)$. Write the Lefschetz decomposition of $u$ as

\begin{equation}
    \begin{aligned}
        u = f \omega^{m-1} + \psi_1 \wedge \omega^{m-2} + \sum\limits_{k =2}^{m-1}\psi_{k} \wedge \omega^{m-1-k}  
    \end{aligned}
\end{equation}


\noindent where $f \in C^{\infty}(M)$ and $\psi_{k} \in C^{\infty}_{k,k}(M)$ are the primitive coordinates of $u$ that satisfy $\Lambda(\psi_k) = 0$, or equivalently $\psi_{k} \wedge \omega^{n-2k+1} = 0$ for each $1 \leq k \leq m-1$. Then the Kernel condition becomes a system of first-order equations for primitives as follows. K\"ahler identity $[L, \sbpartial] = -i \partial $ implies the formula $\sbpartial (v \wedge \omega^p) = \sbpartial v \wedge \omega^{p}+ p \;i \partial v \wedge \omega^{p-1}$, so that

\begin{equation}\label{del-primitive}
    \begin{aligned}
        \sbpartial u = &\sbpartial f \wedge \omega^{m-1} + ( i(m-1)\partial f + \sbpartial \psi_1) \wedge \omega^{m-2} \\
        &+ \sum_{k=2}^{m-1}((m-k)i\partial \psi_{k-1} + \sbpartial \psi_{k} )\wedge\omega^{m-k-1}.
    \end{aligned}
\end{equation}


This is not a primitive decomposition since $\partial \psi_k$ is not a primitive, although $\sbpartial \psi_{k}$ is. Let

\begin{equation}\label{prim-decomp2}
\partial \psi_k = (\partial \psi_k)_{\text{prim}} + \eta_k \wedge \omega 
\end{equation}

\

\noindent be the primitive decomposition of $\partial \psi_k$, for $\eta_k\in C^{\infty}_{k,k-1}(M)$ a primitive form. The fact that $\eta_k$ is primitive follows from the following. The commutation relation $[\Lambda, \partial] = i\bar{\partial}^*$ implies $\Lambda(\partial \psi_k) = i \sbpartial \psi_k$. So $\Lambda(\partial \psi_{k})$ is primitive. If 

$$\eta_k = \eta_k^0 + \eta_k^1 \wedge \omega + \dots $$

\noindent is a primitive decomposition, then applying $\Lambda$ to \eqref{prim-decomp2} shows

\begin{equation}
    \begin{aligned}
    \Lambda(\partial\psi_k) = c_0 \eta_k^0 + c_1 \eta_k^1 \wedge \omega + \dots
    \end{aligned}
\end{equation}


\noindent for non-zero constants $c_i$. This means $\eta_k^i=0$ for $i\geq 2$ and $\eta^0_k$ is primitive.





% Also

% $$\langle \Lambda(\partial \psi_k), \beta\rangle = \langle \psi_k, \spartial \beta \wedge \omega - i \bpartial \beta\rangle = \langle i \sbpartial  \psi_k, \beta\rangle$$


% \

% \noindent for any $\beta$ since $\Lambda(\psi_k) = 0$.


\

We also have

\begin{equation}
    \begin{aligned}
        i \sbpartial \psi_k &= \Lambda(\partial \psi_k) \\
        &= \Lambda(\eta_k \wedge \omega) = (n-2k+1) \eta_k.
    \end{aligned}
\end{equation}

Hence

\begin{equation}
    \begin{aligned}
        (\partial \psi_k)_{\text{prim}} = \partial \psi_k - \frac{i \sbpartial \psi_k \wedge \omega}{n-2k+1}.
    \end{aligned}
\end{equation}

\


Plug this into \eqref{del-primitive} and using $\sbpartial f  =0$ we get the primitive decomposition of $\sbpartial u$.

\begin{equation}
    \begin{aligned}
        \sbpartial u &= (\sbpartial \psi_1 + i(m-1)\partial f + (m-2)i \eta_1) \wedge \omega^{m-2} \\
        &\;\;+ \sum_{k=2}^{m-1}\left[(m-k)i(\partial \psi_{k-1})_{prim} + \sbpartial \psi_{k} + (m-k-1)i\eta_{k}\right]\wedge\omega^{m-k-1}\\
        &=\left[\frac{n-m+1}{n-1}\sbpartial \psi_1 + i(m-1)\partial f \right] \wedge \omega^{m-2}+ \sum_{k=2}^{m-1}\Big[ (m-k)i\partial \psi_{k-1} \\
        &\;\;+ \frac{m-k}{n-2k+ 3} \sbpartial \psi_{k-1} \wedge \omega + \frac{n-m-k+2}{n-2k+1}\sbpartial \psi_{k} \Big]\wedge\omega^{m-k-1}
    \end{aligned}
\end{equation}

The Kernel condition $\sbpartial u =0 $ now gives the following general identity relating primitives of $u$.

\begin{equation}\label{cascade}
    (m-k)i\partial \psi_{k-1} + \frac{m-k}{n-2k+ 3} \sbpartial \psi_{k-1} \wedge \omega + \frac{n-m-k+2}{n-2k+1}\sbpartial \psi_{k} = 0
\end{equation}



\

\noindent for each $1 \leq k \leq m-1 $, assuming $\psi
_0 = f$. For $k=1$, this gives



\begin{equation}\label{primitives12}
    \sbpartial\psi_{1} = -\frac{i(m-1)(n-1)}{n-m+1} \partial f.
\end{equation}

\ 

We note the following implication which will not be used further. 

\begin{lemma}
    If $u$ satisfies $\sbpartial u =0$, then $\partial \sbpartial \psi_{k} = 0$ for all $0 \leq k \leq m-2$ and $k \neq \dfrac{n}{2}$. 

    \end{lemma}

\begin{proof}
    Applying $\sbpartial$ to \eqref{cascade}, we get

\begin{equation}
    \begin{aligned}
        (m-k)i \sbpartial \partial \psi_{k-1} + \frac{(m-k)}{(n-2k+3)} i \partial \sbpartial \psi_{k-1} = 0
    \end{aligned}
\end{equation}

By the anti-commutation $\partial \sbpartial = -\sbpartial \partial$, we get the desired result.



\end{proof}


We recall the Dolbeault Green's operator $G_\partial : C^\infty_{p,q}(M) \to C^\infty_{p,q}(M)$ which is the unique bounded linear operator that commutes with $\Delta_\partial$ and satisfies $I = \mathcal{H}_\partial + \Delta_\partial G_\partial$, where $\mathcal{H}_\partial$ is the orthogonal projection onto the space of harmonic forms. In addition, $G_{\partial}$ satisfies

\

\begin{enumerate}
    \item $G_{\partial}(\alpha) = 0$ for harmonic forms $\alpha$.
    \item $G_{\partial}(\alpha)$ gives the unique solution to $\Delta_{\partial} \beta = \alpha$, that is orthogonal to harmonic forms assuming $\alpha$ is orthogonal to harmonic forms. 

    \item $G_{\partial}$ commutes with $\partial$ and $\spartial$.
\end{enumerate}

\

It is known that if $\partial \alpha = \beta$, $\alpha$ must have the form

\begin{equation}
    \alpha = G_{\partial} \partial^* \beta + \gamma.
\end{equation}

\noindent for any $\partial$-closed form $\gamma$. In fact, this is an immediate consequence of the Hodge theorem which on compact complex manifolds gives the decomposition

$$\alpha = \partial (G_{\partial} \spartial \alpha) + \spartial(G_{\partial} \partial \alpha) + \alpha_{H},$$

\

\noindent for $\alpha_H$ a harmonic form \cite[Chapter $0$]{Griffiths-Harris}. If $\alpha$ is real and $\sbpartial$-exact, the solution is unique with $\gamma =0$. From \eqref{cascade} $\psi_{k}$'s must have the form

\begin{equation}\label{dbar-problem}
    \begin{aligned}
         \psi_{k-1} = G_{\partial} \partial^*\left( \frac{i}{n-2k+ 3} \sbpartial \psi_{k-1} \wedge \omega + \frac{i(n-m-k+2)}{(m-k)(n-2k+1)}\sbpartial \psi_{k} \right) + \gamma
    \end{aligned}
\end{equation}

\noindent for all $2 \leq k \leq m-1$.



\






\begin{claim}\label{harmonicity-claim}
    In the above setting if $\psi_k$ is $\sbpartial$-exact for all $1 \leq k \leq m-2 $ and $\sbpartial \psi_{m-1} = 0$, then all the intermediate primitives $\psi_{k}$ for $1 \leq k \leq m-2$ are zero, and $f$ is a constant. If $\psi_k$'s are assumed to be only $\sbpartial$-closed, then all primitives up to $k=m-2$ are harmonic.
\end{claim}




\
% ********************************

% \begin{proof}

% Since $\sbpartial \psi_{m-1} = 0$, we have that 


% \begin{equation}
%     \begin{aligned}
%         0&= i\partial \psi_{m-2} + \frac{1}{n-2m+ 5} \sbpartial \psi_{m-2} \wedge \omega  \\
%         & = i (\partial \psi_{m-2})_{prim} +\left(\frac{1}{n-2m+ 5} - \frac{1}{n-2m+3}\right)\sbpartial \psi_{m-2} \wedge \omega 
%     \end{aligned}
% \end{equation}


% By uniqueness of primitive decomposition, it follows that 

% $$\sbpartial \psi_{m-2} = 0, \text{  and  }\;\;\; \partial (\psi_{m-2})_{prim} = 0.$$


%  This implies that $\partial \psi_{m-2} = 0$, and hence $\Delta \psi_{m-2} =0$. 

    
% \end{proof}


% *********************************
\begin{proof}


From \eqref{dbar-problem}, it follows that 

\begin{equation}\label{claim1.1-1}
    \psi_{m-2} = G_{\partial} \partial^*\left( \frac{i}{n-2m+ 5} \sbpartial \psi_{m-2} \wedge \omega \right)
\end{equation}

Let $D_k$ be the differential operator 

$$D_k \alpha =\frac{i}{n-2k+3} G_{\partial} \partial^*\left(  L \sbpartial \alpha \right).$$

\ 

Then we show that $D_k(\alpha) = \alpha$ can have only the trivial real-solution. Let $\alpha$ be such a solution. Since $G_{\partial}$ commutes with $\partial^*$, $\alpha$ is $\partial^*$-exact. This means $\alpha$ is orthogonal to harmonic forms by the Hodge theorem. By K\"ahler commutation relations

\begin{equation*}
\begin{aligned}
       D_k(\alpha) &= \frac{i}{n-2k+3} G_{\partial}(L \partial^* \sbpartial \alpha - i \bpartial \sbpartial \alpha) \\
       &= \frac{i}{n-2k+3} \left(-G_{\partial}(L \sbpartial \partial^* \alpha) - i \bpartial G_{\partial}(\sbpartial \alpha)\right) = \frac{1}{n-2k+3}  \bpartial G_{\partial}(\sbpartial \alpha)
\end{aligned}
\end{equation*}


\noindent since $\alpha$ is $\partial^*$-exact and $G_{\partial}$ commutes with $\bpartial$. It follows that $\alpha$ must be $\bpartial$-exact, and


$$\bpartial \alpha = \bpartial D_k(\alpha) = 0.$$

But since $\alpha$ is real, this implies $\partial \alpha =0$. So 

$$\Delta \alpha = (\partial \partial^* +\partial^* \partial) \alpha = 0. $$

Hence it follows that $\alpha = 0$, since it is also orthogonal to harmonic forms. 

\



Using this in \eqref{claim1.1-1} and $\sbpartial \psi_{m-1} =0$, we infer that $\psi_{m-2} = 0$. By induction, it follows that $\psi_{k} =0$ for all $1 \leq  k \leq m-2$. Clearly, if $\psi_1 = 0$, $f$ is a constant by \eqref{primitives12}. The second part of the Claim follows directly from  \eqref{cascade}.



% ***************
% From \eqref{prim-decomp2}

% \begin{equation}
%     \begin{aligned}
%      |\partial \psi_{i}|_{\omega}^2 &= |\partial (\psi_i)_{\text{prim}}|^2_{\omega} + |\eta_i \wedge \omega |_{\omega}^2\\
%      &= |\partial (\psi_i)_{\text{prim}}|^2_{\omega} + c_{n,m,i}|\sbpartial \psi_i|_{\omega}^2
%     \end{aligned}
% \end{equation}

% Since from \eqref{cascade}, $(m-i)i(\partial \psi_{i-1})_{\text{prim}} + \sbpartial \psi_{i} =0$ for all $i$, it follows that if $\sbpartial \psi_{i} =0$ then $(\partial \psi_{i-1})_{\text{prim}}=0$ for all $i$, and

% \begin{equation}
%     \begin{aligned}
%      |\partial \psi_{i-1}|_{\omega}^2&=  c_{n,m,i-1}|\sbpartial \psi_{i-1}|_{\omega}^2
%     \end{aligned}
% \end{equation}

% Now from \eqref{cascade}, if $\sbpartial \psi_{i} =0$ then

% $$\sbpartial \psi_{i-1} \wedge \omega =0.$$

% This implies $\sbpartial \psi_{i-1} = 0$, since $[\Lambda, \sbpartial] = 0$ implies $\sbpartial \psi$ is a primitive. Now by induction it follows that if $\sbpartial \psi_{m-1} = 0$, then $\partial \psi_{i}=0$ and  $\sbpartial \psi_{i} =0$ for all $1 \leq i \leq m-2$. This implies harmonicity of $\psi_i$'s and $\partial f= 0$. So $f$ is constant.


   
\end{proof}




\


% Hodge decomposition says that $\psi_1 = \bpartial \alpha + \sbpartial \beta + \gamma$, where $\alpha \in C^{\infty}_{1,0}(M)$, $\beta \in C^{\infty}_{1,2}(M)$ and $\gamma$ is an $\omega$-harmonic $(1,1)$-form. From this and \eqref{primitives12}

% \begin{equation}
%     \sbpartial \psi_1 = \sbpartial \bpartial \alpha = -\frac{i(m-1)(n-1)}{n-m+1} \partial f
% \end{equation}

% \

% Using this we evaluate $\Delta \psi_1$.

% \begin{equation}
%     \begin{aligned}
%         \Delta \psi_1 & = \bpartial \sbpartial \psi_1+ \sbpartial\bpartial \psi_1\\
%         & = \bpartial \sbpartial \bpartial \alpha + \sbpartial \bpartial \sbpartial \beta \\
%         & = \frac{(m-1)(n-1)}{n-m+1} i\partial \bpartial f +  \Delta \sbpartial \beta 
%     \end{aligned}
% \end{equation}


\begin{corollary}\label{lemma3}
    Let $u^1$ and $u^2$ be two real $(m-1,m-1)$ forms in $\Ker (\spartial) \cap \Ker( \bar{\partial}^*)$. Then assuming $\psi^1_{k} -\psi^2_{k}$ is $\sbpartial$-exact for all $1 \leq k \leq m-2$ and $\psi^1_{m-1} -\psi^2_{m-1}$ is $\sbpartial$-closed 
    we have 
    
    $$u^1(z) = u^2(z) + c \omega^{m-1} + (\psi^1_{m-1} -\psi^2_{m-1}).$$
    
\end{corollary}


\ 



Define

\begin{equation}
Q(u) = \star_{\omega}(i \pbp u \wedge \omega_{n-m-1}).
\end{equation}

\

Then \cite[Cor.~3.4]{DP25} shows that

\begin{equation}\label{defineQ}
        Q(u)= \frac{1}{(m-1)!}\left[-i \pbp \Lambda^{m-1} u + (m-1) \Delta \Lambda^{m-2} u - (\Delta \Lambda^{m-1} u) \;\omega\right].
\end{equation}



\

So the Monge-Amp\`ere equation \eqref{MA-eqn} can now be written as 

\begin{equation}\label{QMA-eqn}
    \left[\star_{\omega}(\alpha\wedge \omega_{n-m-1}) + Q(u)\right]^n = dV.
\end{equation}

\



By taking the Trace of $u$ repeatedly it can be shown that 

\begin{equation}\label{trace-2}
    \Lambda^{m-2} u =\frac{(m-2)!(n-2)!}{(n-m)!} \psi_1 + \frac{(m-1)!(n-1)!}{(n-m+1)!} f \; \omega,
\end{equation}

\noindent and 

\begin{equation}\label{trace-1}
\Lambda^{m-1} u = \frac{(m-1)! n!}{(n-m+1)!}f.
\end{equation}

Using \cref{defineQ,trace-1,trace-2}, equation \eqref{QMA-eqn} can be written as

\begin{equation}\label{QMA-primitives}
    \Big[\star_{\omega}(\alpha\wedge \omega_{n-m-1}) + \frac{(n-2)!}{(n-m)!}\left(\Delta \psi_1 - c_{n,m}\left[(n-m+1)\Delta f \omega +  n i \pbp f \right]\right)\Big]^n = dV.
\end{equation}

\noindent where $c_{n,m} =\dfrac{(n-1)}{(n-m+1)}$. 


\

We prove a uniqueness theorem for \eqref{MA-eqn}.








\begin{theorem}\label{unique-thm}
Let $u^1$ and $u^2$ in $C_{m-1,m-1}^{\infty}(M, \mathbb R)$ be two solutions to \eqref{MA-eqn} that lie in $\Ker (\spartial) \cap \Ker( \bar{\partial}^*)$. Then

\begin{enumerate}
    \item If the difference of all primitives of orders from $2$ to $(m-2)$ of $u^1$ and $u^2$ is $\bar\partial^*$-exact and $\sbpartial(\psi^1_{m-1} - \psi_{m-1}^2) =0 $, then

$$u^1(z) = u^2(z) + c \omega^{m-1} + \alpha \wedge \omega^{m-2} +  (\psi^1_{m-1} - \psi_{m-1}^2).$$

\noindent for some constant $c$, a harmonic $(1,1)$ form $\alpha$.

\

\item If the difference of all primitives of orders from $2$ to $(m-1)$ of $u^1$ and $u^2$ is $\bar\partial^*$-closed, then

$$u^1(z) = u^2(z) + c \omega^{m-1} + \sum_{k=1}^{m-2} \alpha_k \wedge \omega^{m-k-1} + (\psi^1_{m-1} - \psi_{m-1}^2).$$

\noindent for some constant $c$ and harmonic forms $\alpha_k \in C^{\infty}_{k,k}(M, \mathbb R) $.
\end{enumerate}


\end{theorem}

\

\begin{proof}


We write $\Delta \psi_1$ in terms of $f$ and $\psi_2$ as follows. First apply $\spartial$ to \eqref{cascade} with $k=2$ to get (assuming $m\neq 2$)

\begin{equation}
    \begin{aligned}
        \spartial \partial \psi_{1} = \frac{1}{(n-1)} \left(\bpartial \sbpartial \psi_1 + i \spartial \sbpartial \psi_1 \; \omega\right) + i \frac{n-m}{(n-3)(m-2)} \spartial \sbpartial \psi_2 
    \end{aligned}
\end{equation}

It follows from \eqref{primitives12} that 

$$ \bpartial \sbpartial \psi_1 = -\frac{(m-1)(n-1)}{n-m+1} i\bpartial \partial f$$

\noindent and

$$i \spartial \sbpartial \psi_1 \; \omega = \frac{(m-1)(n-1)}{n-m+1} \Delta f \; \omega.$$


Combining the above equations we get

\begin{equation}
\begin{aligned}
    \Delta \psi_1 &= \partial \spartial \psi_1 + \spartial \partial \psi_1\\
    &= \frac{(m-1)n}{n-m+1} i \partial \bpartial f + \frac{m-1}{n-m+1} \Delta f \; \omega + \frac{i(n-m)}{(n-3)(m-2)} \spartial \sbpartial \psi_2.
\end{aligned}
\end{equation}

Plugging this into \eqref{defineQ} using \eqref{trace-2} and \eqref{trace-1} gives the following.


\begin{equation}\label{Q-fp1}
    Q(u) = \frac{(n-2)!}{(n-m)!} \left[ \frac{n(m-n)}{n-m+1} (\Delta f \omega +   i \partial \bar{\partial} f) + i \frac{n-m}{(m-2)(n-3)} \partial^* \bar{\partial}^* \psi_2 \right]
\end{equation}


For $m=2$, the term $\partial^* \bar{\partial}^* \psi_2$ will be absent in the above formula. Now the PDE is expressed in terms of $f$ and the primitive $\psi_2$. If $u_1$ and $u_2$ are two solutions, then $\Lambda_{\rho}(Q(u_1 - u_2)) = 0$ with $\rho >0$ is a $(1,1)$ form defined by

$$\rho^{n-1} = \sum_{s=1}^{n} \left[ \star_{\omega} \Big( (\alpha + i\partial\bar{\partial}u_1) \wedge \omega_{n-m-1} \Big) \right]^{n-s} \wedge \left[ \star_{\omega} \Big( (\alpha + i\partial\bar{\partial}u_2) \wedge \omega_{n-m-1} \Big) \right]^{s-1},$$


\noindent which follows from the binomial expansion of

\begin{equation}
    \begin{aligned}
         (\star_{\omega}\left((\alpha + i \pbp u_1)\wedge \omega_{n-m-1}\right))^n -  (\star_{\omega}\left((\alpha + i \pbp u_2)\wedge \omega_{n-m-1}\right))^n = 0.
    \end{aligned}
\end{equation}




 By \eqref{Q-fp1} this implies that

\begin{equation}\label{max-principle}
    \begin{aligned}
    0 &=\Lambda_{\rho}\left(\frac{n(m-n)}{n-m+1} \Delta (f_1 -f_2) \; \omega +   \frac{n(m-n)}{n-m+1} i \partial \bar{\partial} (f_1 - f_2)\right),
    \end{aligned}
\end{equation}

\noindent since $\psi^1_{2}-\psi^2_{2}$ is $\bpartial^*$-closed by assumption. Define the operator

$$\mathcal L v = \frac{n(m-n)}{n-m+1}\left( \mbox{Tr}_{\rho}(\omega)  \; \omega^{i \bj}  +  \rho^{i \bj}\right)v_{i \bj}.$$


This is an elliptic operator. By \eqref{max-principle} $\mathcal L(f_1 - f_2) = 0$ on a compact manifold. Hence the maximum principle implies $f_1 - f_2 = \text{constant}$.


Now Theorem \ref{unique-thm} follows from the proof of Claim \eqref{harmonicity-claim} combined with the fact that if $f_1 - f_2$ is constant, then $\spartial(\psi^1_1  - \psi^2_1) = 0$. From equation \eqref{cascade}, we also get that $\partial(\psi^1_1  - \psi^2_1) = 0$. Hence $\psi^1_1  - \psi^2_1$ is harmonic. 



\end{proof}


% \begin{remark}
%     It is possible to write $Q(u)$ in terms of $f$ and the highest primitive $\psi_{m-1}$ using \eqref{cascade}. But the final formula will be quite lengthy. 
% \end{remark}



% We denote the primitive coordinates of $u^1$ and $u^2$ by $(f^1,\psi_i^1)$ and $(f^2, \psi_i^2)$ respectively. Then by assumption and Claim \ref{harmonicity-claim}, $f^1-f^2$ is constant and $\psi_i^1 - \psi_i^2$ is harmonic. So

% $$\Lambda^{m-2} u$$

% By \eqref{QMA-eqn} 

% \begin{equation}
%     \begin{aligned}
%     0 &= \Lambda_{\rho} (Q(u^1-u^2))\\
%     & = \Lambda_{\rho}()
%     \end{aligned}
% \end{equation}
% (expand on this more.)

{\flushleft{\bf The Laplacian Trace condition:} }

\bigskip

In addition to above, a third condition is introduced in \cite{DP25}:

\

\begin{equation}\label{MA-eqn2}
\begin{aligned}
    &\left[\star_{\omega}\left((\alpha + i \pbp u)\wedge \omega_{n-m-1}\right)\right]^n = dV\\
\end{aligned}
\end{equation}


\noindent with $u$ satisfying 

\begin{equation}\label{generalIC}
    \alpha + i \pbp u >0, \quad u \in \Ker (\spartial) \cap \Ker( \bar{\partial}^*), \quad \text{ and }\Lambda^{m-2} \Delta u =0.
\end{equation}

\

\ This third condition is too restrictive since along with \eqref{trace-2} it implies


\begin{equation}
    \frac{(m-2)!(n-2)!}{(n-m)!} \Delta \psi_1 + \frac{(m-1)!(n-1)!}{(n-m+1)!} \Delta f \omega = 0.
\end{equation}

\


Here we used that $\Delta$ commutes with $\Lambda$. This commutation also implies that $\Delta \psi_1$ and $\Delta f$ are both primitives. So by the uniqueness of Lefschetz decomposition, $\Delta \psi_1 = \Delta f =0$. This means that $f$ is constant. So now equation \eqref{QMA-primitives} reads

$$ [\star_{\omega}(\alpha\wedge \omega_{n-m-1})]^n = dV,$$


\

\noindent which cannot hold for arbitrary $dV$. In view of this, we introduce an alternate gauge-fixing in the next section.


% \begin{proposition}
%     The system \eqref{MA-eqn2} reduces to the complex Monge-Amp\`ere equation for $f$:

%     \begin{equation} \label{QMA-eqn2}
%         \left(\Omega + \frac{n!}{(n-m+1)!} \;i \pbp f\right)^n = dV,
%     \end{equation}

%     \

%     \noindent where $\Omega = \star_{\omega} (\alpha\wedge \omega_{n-m-1})$ with 
%     $$\Omega + \frac{n!}{(n-m+1)!} \;i \pbp f>0.$$
% \end{proposition}


% \

% \begin{proof}

% Since $\Delta$ commutes with the Trace operator $\Lambda$ and since $\Delta$ preserves primitives, $\Lambda^{m-2} \Delta u =0$ implies that

% $$ \frac{1}{(n-m+1)} \Delta \psi_1 + (m-1)(n-1) \Delta f \; \omega = 0,$$

% \noindent from \eqref{trace-2}. This is a primitive decomposition and so we have that $\Delta \psi_1 = \Delta f =0$. Then from \eqref{QMA-primitives} 

% \begin{equation}
%     Q(u) = \frac{n!}{(n-m+1)!} \;i \pbp f.
% \end{equation}
% \end{proof}



\section{$(a,b)$ Monge-Amp\`ere-type equation}\label{section3}

\bigskip


 Fix any two real constants $(a,b)$ satisfying $a+ nb = m-n-2$. Consider the equation

\begin{equation} \label{(n-1)-MAeqn}
    \begin{aligned}
        \left[\star_{\omega}(\alpha + i\pbp u) \wedge \omega_{n-m-1}\right]^n = dV
    \end{aligned}
\end{equation}

\


\noindent with conditions $\alpha + i\pbp u >0$, $u \in \Ker (\spartial) \cap \Ker( \bar{\partial}^*)$, and 

\

\begin{equation}\label{alt-third-condition}
     (m-1)\Lambda^{m-2} \Delta u = (1+a) i \pbp \Lambda^{m-1}  u +  (1+b) \Lambda^{m-1} \Delta u \;\omega.
\end{equation}


\

Taking $\Lambda$ on both sides of \eqref{alt-third-condition} shows that $a+nb = m-n-2$ is necessary. Under \eqref{alt-third-condition}, equation \eqref{(n-1)-MAeqn} simplifies to 

\begin{equation}\label{ab-MAeqn}
    \left(\star_{\omega}(\alpha \wedge \omega_{n-m-1}) + \frac{n!}{(n-m+1)!} \left[ a i\pbp f + b \Delta f \omega\right] \right)^n = dV.
    \end{equation}

Indeed, $Q(u)$ now becomes

\begin{equation}
    \begin{aligned}
        Q(u) &= \frac{1}{(m-1)!}\left(a  i \pbp \Lambda^{m-1}  u +  b \Lambda^{m-1} \Delta u \;\omega\right)\\
    \end{aligned}
\end{equation}

 Now using \eqref{trace-1} we get \eqref{ab-MAeqn}. We will call \eqref{ab-MAeqn} the {\em $(a,b)$ Monge-Amp\`ere-type equation}.

\

We show that the application to Demailly's inequality holds, using the $(a,b)$ Monge-Amp\`ere-type equation for $b=0$ and $a= m-n-2$. We adapt the analytic technique developed in \cite{DP25} and \cite{Popovici16}.


\begin{theorem}\label{demailly-thm}
    Let $m \in \{1, \hdots, n-1\}$. Suppose $\alpha, \beta \in C^{\infty}_{m,m}(M, \mathbb R)$ are such that

    \begin{enumerate}
        \item $d\alpha =0 $, $\alpha >0$ and $\Delta \alpha_{\omega} =0$.

        \ 

        \item $\beta \geq C \omega_m$ for some constant $C>0$ and $\pbp \beta =0$.

        \ 

        \item $\dfrac{1}{n!}\int_M \alpha_{\omega}^n - \dfrac{1}{(n-m)!}\int_M\alpha_{\omega}^{n-m}\wedge \beta >0$
    \end{enumerate}

    \

    Then there exists a real $(m,m)$ current $T$ such that 

    $$T \geq \delta (\alpha_{\omega})_m \text{ for some constant } \delta >0,\hspace{0.5cm} \pbp T = 0, \text{ and } \hspace{0.5cm} T \in [(\alpha_{\omega})_m - \beta]_A.$$

    If in addition, $d \beta = 0$, $T$ can be found such that $dT = 0$ and $T$ lies in $[(\alpha_{\omega})_m - \beta]_{BC}$. Here $[.]_{BC}$ stands for the Bott-Chern cohomology and $[.]_A$ stands for the Aeppli cohomology class.
\end{theorem}

\begin{proof}

By Lamari-type duality lemma \cite{Lamari99, DP25}, to prove that there exists a real current $S$ and a constant $\delta>0$ such that $T = (\alpha_{\omega})_{m} - \beta + i \pbp S$ satisfies the Theorem, it is enough to show that there exists a constant $\delta>0$ such that 

$$(1-\delta)\int_{M}(\alpha_{\omega})_{m}\wedge\Omega\geq\int_{M}\beta\wedge\Omega,$$

  \noindent for any $\Omega \in C^{\infty}_{n-m,n-m}(M, \mathbb R)$ satisfying $\pbp \Omega = 0$ and $\Omega>0$ weakly. Assume for contradiction that there is a sequence of forms $\Omega_k$ and $\delta_k \to 0$ such that 

  
  \begin{equation}\label{contra-sequence}
       (1-\delta_k)\int_{M}(\alpha_{\omega})_{m}\wedge\Omega_k<\int_{M}\beta\wedge\Omega_k.
  \end{equation}
 
  By assumption $2$, it is possible to normalize the sequence $\Omega_k$ so that $\int_{M}\beta\wedge\Omega_k = 1$ for each $k$. This also implies $$\int_M \omega_m \wedge \Omega_k \le \frac{1}{C},$$

  \noindent so that $\Omega_k \to \Omega_{\infty}$ weakly. Taking limits of \eqref{contra-sequence} gives
  
  $$\lim_{k \to \infty} \int_{M} (\alpha_{\omega})_{m} \wedge \Omega_{k} = \int_{M} (\alpha_{\omega})_{m} \wedge \Omega_{\infty} \leq 1.$$
  
  Let $u_k$ be a solution of 

  \begin{equation}
      \frac{1}{n!} \left[ \star_\omega \Big( (\alpha + i\partial\bar{\partial}u_k) \wedge \omega_{n-m-1} \Big) \right]^n = c_k \beta \wedge \Omega_k
  \end{equation}

  \noindent such that 

  \begin{equation}
      \tilde{\alpha}_k = \alpha + i\partial\bar{\partial}u_k > 0 \quad \text{and} \quad u_k \in \Ker(\partial^*) \cap \Ker(\bar{\partial}^*),
  \end{equation}

  \noindent and

  \begin{equation}\label{generalIC1}
      (m-1)\Lambda^{m-2} \Delta u_k = (m-n-1)i \pbp \Lambda^{m-1}  u_k +  \Lambda^{m-1} \Delta u_k \;\omega.
  \end{equation}   

  \

  \noindent which exists by Theorem \ref{exist-theorem}.

  \

  By H\"older's inequality and the pointwise inequality $\dfrac{\rho_m \wedge \Omega}{\omega_n} \cdot \dfrac{\rho_{n-m} \wedge \beta}{\rho_n} \geq \dfrac{\beta \wedge \Omega}{\omega_n}$ for any $(1,1)$ form $\rho>0$ (See \cite[p.~25]{DP25})

$$\begin{aligned}
&\left( \int_M ((\tilde{\alpha}_k)_\omega)_m \wedge \Omega_k \right) \cdot \left( \int_M ((\tilde{\alpha}_k)_\omega)_{n-m} \wedge \beta \right) \geq c_k \left( \int_M \beta \wedge \Omega_k \right)^2 = c_k.
\end{aligned}$$


Note that $c_k = \int_M ((\tilde{\alpha}_k)_\omega)_n $. By \eqref{generalIC1}

$$\begin{aligned}
(\tilde{\alpha}_k)_\omega &= \star_{\omega}\left(\alpha \wedge \omega_{n-m-1}\right) + Q(u_k) \\
&= \alpha_\omega  + \frac{m-n-2}{(m-1)!}\left( i \pbp \Lambda^{m-1}  u_k \right).
\end{aligned}$$

\

It follows that for any $p \in \{1,\hdots, n\}$

$$((\tilde{\alpha}_k)_\omega)_p - (\alpha_\omega)_p = i\pbp \Lambda^{m-1} u_k \wedge \Gamma ,$$

\

\noindent for a $d$-closed form $\Gamma$. Note that $\alpha_{\omega}$ is harmonic and hence $d$-closed. Now applications of the Stokes' theorem using $\pbp \Omega_k=0$ and $\pbp \beta = 0$ gives

$$ \int_M ((\tilde{\alpha}_k)_\omega)_m \wedge \Omega_k = \int_M (\alpha_\omega)_m \wedge \Omega_k,$$

$$ \int_M ((\tilde{\alpha}_k)_\omega)_{n-m} \wedge \beta =  \int_M (\alpha_\omega)_{n-m} \wedge \beta , $$


\noindent and

$$\int_M ((\tilde{\alpha}_k)_\omega)_n =\int_M ({\alpha}_\omega)_n .$$


\

Combining the above equations with $\int_{M} (\alpha_{\omega})_{m} \wedge \Omega_{\infty} \leq 1$, we have


$$\int_{M} (\alpha_{\omega})_{n-m} \wedge \beta \geq \int_{M} (\alpha_{\omega})_{n}$$

\

\noindent contradicting the assumption $3$.



\end{proof}


\begin{thebibliography}{99}
\bibitem[]{DP25} Dinew, S\l awomir; Dan Popovici. m-Pseudo-effectivity and a Monge-Ampère-Type Equation for Forms of Positive Degree. arXiv:2510.27362 (2025)
\bibitem[]{Griffiths-Harris} Griffiths, Phillip; Harris, Joseph. Principles of algebraic geometry. (1994) xiv+813
\bibitem[]{Lamari99} Lamari, Ahc\`ene. Courants k\"ahl\'eriens et surfaces compactes. Ann. Inst. Fourier (Grenoble) 49 (1999) vii, x, 263--285
\bibitem[]{Popovici16} Popovici, Dan. Sufficient bigness criterion for differences of two nef classes. Math. Ann. 364 (2016) 649--655
\end{thebibliography}
\end{document}
